\section{Introduction}

The arrangement graph $A(n,k)$ has as vertices the injective words of length
$k$ over an alphabet of size $n$; two vertices are adjacent when they differ
in exactly one coordinate. It is a standard model for interconnection
networks. We study the fixed-volume external vertex-boundary problem: for a
given $R$, how small can $|N(S)|$ be among $R$-vertex subsets $S$?

This problem is related to, but distinct from, $g$-extraconnectivity
$\kappa_g(G)$, the minimum size of a vertex cut whose removal leaves every
component larger than $g$. A small-boundary set supplies a candidate cut only
when deleting its external boundary has the required component sizes; a
general converse reduction is also needed. We do not claim that reduction or
an all-parameter formula for $\kappa_g(A(n,k))$ here. Existing work on
arrangement-graph extra-connectivity and fault diagnosis provides context for
these questions~\cite{cheng2022extraconnectivity,cheng2013linearly}.

Our main unconditional result is a Hamming comparison with an explicit
volume-only additive error. The coordinate-root count gives the Hamming
linear slope, while a global overlap-charging argument controls the loss from
multiple coordinate descriptions of the same external neighbor. When a
Boolean Hamming ball embeds, it provides a witness attaining the comparison
value. Thus the minimum boundary lies in a proved interval, but need not equal
the Hamming value.

The distinction is essential. A full radius-one Star in $A(11,6)$ has
volume $31$ and boundary $450$, while the embedded Hamming-ball comparison is
$476$. This refutes the former exact-minimum claim even under its proposed
embedding gate. The additive allowance in our theorem is therefore not a
technical artifact that can be deleted without further hypotheses.

The paper develops the following components:
\begin{itemize}
\item We formalize the coordinate-root boundary identity and prove the
  hypercube-derived defect bound for arbitrary subsets.
\item We prove, using Bonferroni and a global witness-pair injection, a
  dimension-independent collision bound and the resulting unconditional
  Hamming sandwich (Section~\ref{sec:hamming-sandwich}).
\item We retain the fixed-topology penalty calculation as an asymptotic
  comparison tool, while distinguishing it from a classification of all
  extremizers or a uniform phase diagram.
\end{itemize}

The Lean development checks the boundary identity, defect estimate, collision
charging, and sandwich theorem. The remaining extremal question is to
characterize the lower envelope of Hamming-like, Star-like, and intermediate
configurations, especially at volumes comparable to $k(n-k)$. Connecting
that profile to extra-connectivity is a separate open task. The underlying
coordinate-root viewpoint is inspired by classical edge-isoperimetric
methods, including Harper's theorem for the cube~\cite{harper1966optimal}.
